%------------------------------------------------------------------------------%
\documentclass[fleqn,utf8,aspectratio=169,12pt]{beamer}

\usepackage{gaal}

\title{Diagonalização de Matrizes Simétricas}

%------------------------------------------------------------------------------%
\begin{document}

%\addtitle
%
%%------------------------------------------------------------------------------%
%\section{Matrizes Simétricas e Matizes Ortogonais}
%
%%------------------------------------------------------------------------------%
%\begin{frame}{Matrizes simétricas}
%  Uma matriz quadrada $A$ é \emph{simétrica} quando
%  \[
%      A^t = A
%  \]
%
%  \pause
%  Exemplo
%  \[
%    A =
%    \begin{pmatrix}
%        2 & 1 & 0\\
%        1 & 2 & 0\\
%        0 & 0 & 4
%    \end{pmatrix}
%  \]
%\end{frame}
%
%%------------------------------------------------------------------------------%
%\begin{frame}{Diagonalização}
%  $A$ é diagonalizável se existe $P$ invertível e $D$ diagonal,
%  \pause
%  \[
%    A = P D P^{-1}
%  \]
%
%  \pause
%  \bigskip
%  As colunas de $P$ são autovetores de $A$
%
%  \pause
%  \bigskip
%  Os elementos da diagonal de $D$ são os autovalores correspondentes
%\end{frame}
%
%%------------------------------------------------------------------------------%
%\begin{frame}{Matriz Ortogonal}
%  Uma matriz quadrada $P$ é \emph{ortogonal} quando
%  \[
%    P^t P = I
%  \]
%
%  \pause
%  isto é
%  \[
%    P^{-1} = P^t
%  \]
%\end{frame}
%
%%------------------------------------------------------------------------------%
%\begin{frame}{Colunas de uma Matriz Ortogonal}
%  \[
%    P =
%    \begin{bmatrix}[c]
%      \; \vec{v}_{1} & \vec{v}_{2} & \cdots & \vec{v}_{n} \;
%    \end{bmatrix}
%  \]
%
%  \[
%    \pause
%    P^t P
%    \pause
%    =
%    \begin{bmatrix}[c]
%      \;\vec{v}^{\,t}_{1}\; \\
%      \vec{v}^{\,t}_{2} \\
%      \vdots  \\
%      \vec{v}^{\,t}_{n}
%    \end{bmatrix}
%    \begin{bmatrix}[c]
%      \; \vec{v}_{1} & \vec{v}_{2} & \cdots & \vec{v}_{n} \; \\
%    \end{bmatrix}
%    \pause
%    =
%    \begin{bmatrix}[c]
%      \; \vec{v}^{\,t}_{1}\vec{v}_{1} & \vec{v}^{\,t}_{1}\vec{v}_{2} & \cdots & \vec{v}^{\,t}_{1}\vec{v}_{n} \; \\[2mm]
%      \; \vec{v}^{\,t}_{2}\vec{v}_{1} & \vec{v}^{\,t}_{2}\vec{v}_{2} & \cdots & \vec{v}^{\,t}_{2}\vec{v}_{n} \; \\[2mm]
%      \vdots                          & \vdots                       & \ddots & \vdots                          \\[2mm]
%      \; \vec{v}^{\,t}_{n}\vec{v}_{1} & \vec{v}^{\,t}_{n}\vec{v}_{2} & \cdots & \vec{v}^{\,t}_{n}\vec{v}_{n} \;
%    \end{bmatrix}
%    \pause
%    = I
%  \]
%\end{frame}
%
%%------------------------------------------------------------------------------%
%\begin{frame}{Colunas de uma Matriz Ortogonal}
%  \[
%    \vec{v}^{\,t}_{i}\vec{v}_{j}
%    \pause
%    \;=\; \delta_{ij}
%    \pause
%    \;=\; \begin{cases}
%      1, & i = j\\
%      0, & i \neq j
%    \end{cases}
%  \]
%
%  \pause
%  As colunas de $P$ são vetores ortonormais
%\end{frame}
%
%%------------------------------------------------------------------------------%
%\section{Diagonalização de Matrizes Simétricas}
%
%%------------------------------------------------------------------------------%
%\begin{frame}{Diagonalização Ortogonal}
%  Uma matriz $A$ é \emph{ortogonalmente diagonalizável}
%
%  \pause
%  quando existem uma \emph{matriz ortogonal} $P$ e uma matriz diagonal $D$
%
%  \pause
%  tais que
%  \[
%    A = PDP^t
%  \]
%\end{frame}
%
%%------------------------------------------------------------------------------%
%\begin{frame}{Teorema Espectral para Matrizes Simétricas}
%  Toda matriz real simétrica é ortogonalmente diagonalizável
%
%  \pause
%  \bigskip
%  Existe uma matriz ortogonal $P$ e uma matriz diagonal $D$ tais que
%  \[
%    A = PDP^t
%  \]
%\end{frame}
%
%%------------------------------------------------------------------------------%
%\begin{frame}{Propriedades das Matrizes Simétricas}
%  \begin{itemize}[<+->]
%    \setlength{\itemsep}{2em}
%    \item Todos os autovalores de $A$ são reais
%    \item Autovetores associados a autovalores distintos são ortogonais
%    \item É possível construir uma base ortonormal de autovetores
%  \end{itemize}
%\end{frame}
%
%%------------------------------------------------------------------------------%
%\begin{proofframe}{Autovetores de Autovalores Distintos}
%  Verificando que autovetores associados a autovalores distintos são ortogonais
%
%  \pause
%  \bigskip
%  Seja $A$ uma matriz simétrica e $\vec{v}$ e $\vec{w}$ autovetores
%
%  \pause
%  associados aos autovalores distintos $\lambda$ e $\mu$
%\end{proofframe}
%
%%------------------------------------------------------------------------------%
%\begin{proofframe}{Autovetores de autovalores distintos}
%  \[
%    \onslide<+->{A\vec{v} = \lambda\vec{v}}  \qquad\qquad
%    \onslide<+->{A\vec{w} = \mu\vec{w}}      \qquad\qquad
%    \onslide<+->{\lambda  \neq  \mu}
%  \]
%  \vspace{-1.5\baselineskip}
%\begin{columns}[t]
%\column{0.25\textwidth}
%  \begin{align*}
%    \onslide<+->{\vec{v}^{\,t} A \vec{w}}
%    \onslide<+->{& = \vec{v}^{\,t} (A \vec{w}) \\}
%    \onslide<+->{& = \vec{v}^{\,t}(\mu\vec{w}) \\}
%    \onslide<+->{& = \mu\vec{v}^{\,t}\vec{w}}
%  \end{align*}
%\column{0.25\textwidth}
%  \begin{align*}
%    \onslide<+->{\vec{v}^{\,t} A \vec{w}}
%    \onslide<+->{& = (\vec{v}^{\,t} A )   \vec{w} \\}
%    \onslide<+->{& = (\vec{v}^{\,t} A^t ) \vec{w} \\}
%    \onslide<+->{& = (A \vec{v})^t    \vec{w} \\}
%    \onslide<+->{& = \lambda\vec{v}^{\,t} \vec{w}}
%  \end{align*}
%\column{0.4\textwidth}
%  \begin{align*}
%    \onslide<+->{\lambda\vec{v}^{\,t}\vec{w}       & = \mu\vec{v}^{\,t}\vec{w} \\}
%    \onslide<+->{\lambda\vec{v}^{\,t}\vec{w} - \mu\vec{v}^{\,t}\vec{w} & = 0   \\}
%    \onslide<+->{(\lambda-\mu)\vec{v}^{\,t}\vec{w} & = 0                       \\}
%    \onslide<+->{\vec{v}^{\,t}\vec{w}              & = 0                       \\}
%    \onslide<+->{\vec{v}                           & \perp \vec{w}}
%  \end{align*}
%\end{columns}
%\end{proofframe}
%
%%------------------------------------------------------------------------------%
%\begin{frame}{Passo a Passo}
%  \onslide<+->{Para diagonalizar ortogonalmente uma matriz simétrica $A$}
%  \begin{enumerate}[<+->]
%      \item Determine os autovalores de $A$
%      \item Determine uma base de cada autoespaço
%      \item Construa uma base ortonormal de autovetores
%      \item Organize os autovetores nas colunas de $P$
%      \item Organize os autovalores correspondentes na diagonal de $D$
%  \end{enumerate}
%  \onslide<+->{Ao final
%  \[
%      A = P D P^t
%  \]}
%\end{frame}
%
%%------------------------------------------------------------------------------%
%\begin{frame}{Autovalores Repetidos?}
%  Uma matriz simétrica pode possuir autovalores repetidos
%
%  \pause
%  \bigskip
%  Nesse caso, o autoespaço associado ao autovalor pode ter dimensão maior que um
%
%  \pause
%  \bigskip
%  É possível escolher uma base ortonormal para cada autoespaço
%
%  \pause
%  \bigskip
%  Se necessário, podemos utilizar o processo de Gram-Schmidt
%\end{frame}
%
%%------------------------------------------------------------------------------%
%\section{Exemplos}
%
%%------------------------------------------------------------------------------%
%%------------------------------------------------------------------------------%
\begin{question}
  Diagonalize ortogonalmente a matriz
  \[
    A =
    \begin{pmatrix}
      5 & 2 \\
      2 & 2
    \end{pmatrix}
  \]
  Determine as matrizes $P$ e $D$ tais que
  \(
    A = P D P^t
  \)
\end{question}

%------------------------------------------------------------------------------%
\begin{resolution}{Autovalores}
\begin{columns}
\column{0.4\textwidth}
  \onslide<+->{Polinômio característico}
  \begin{align*}
    \onslide<+->{p(\lambda)}
    \onslide<+->{& = \det(A-\lambda I)}
    \\
    \onslide<+->{& = \begin{vmatrix}[c]
      5-\lambda & 2         \\
      2         & 2-\lambda
    \end{vmatrix}}
    \\
  \onslide<+->{  & = (5-\lambda)(2-\lambda) - 4 \\}
    \onslide<+->{& = \lambda^{2} - 7\lambda + 6}
  \end{align*}
\column{0.6\textwidth}
  \[
    \onslide<+->{\Delta}
    \onslide<+->{= b^2 - 4ac}
    \onslide<+->{= (-7)^2 - 4 \times 1 \times 6}
    \onslide<+->{= 25}
  \]

  \smallskip
  \[
    \onslide<+->{\sqrt{\Delta}}
    \onslide<+->{= \sqrt{25}}
    \onslide<+->{= 5}
  \]

  \smallskip
  \[
    \onslide<+->{\lambda_1}
    \onslide<+->{= \frac{-b-\sqrt{\Delta}}{2a}}
    \onslide<+->{= \frac{7-5}{2}}
    \onslide<+->{= 1}
  \]

  \smallskip
  \[
    \onslide<+->{\lambda_2}
    \onslide<+->{= \frac{-b+\sqrt{\Delta}}{2a}}
    \onslide<+->{= \frac{7+5}{2}}
    \onslide<+->{= 6}
  \]
\end{columns}
\end{resolution}

%------------------------------------------------------------------------------%
\begin{resolution}{Autovetores para $\lambda=1$}
\begin{columns}
\column{0.4\textwidth}
  \begin{align*}
    \onslide<+->{(A-I)\vec{x}&=\vec{0}}
    \\
    \onslide<+->{\begin{bmatrix}
      4 & 2 \\
      2 & 1
    \end{bmatrix}
    \vetor{x}{y} & = \vetor{0}{0}}
  \end{align*}
  \onslide<+->{\begin{align*}
    4x + 2y & = 0 \\
    2x +  y  & = 0
  \end{align*}}
\column{0.4\textwidth}
  \onslide<+->{Escolhendo $x=1$}
  \onslide<+->{\[
    y = -2
  \]}
  \onslide<+->{\[
    \vec{v}_1 = \vetor{1}{-2}
  \]}
  \onslide<+->{\[
    E(1) = \ger \left\{ \vetor{1}{-2} \right\}
  \]}
\end{columns}
\end{resolution}

%------------------------------------------------------------------------------%
\begin{resolution}{Normalização de $\vec{v}_1$}
  Norma de
  \(
    \vec{v}_1 = \vetor{1}{-2}
  \)
  \pause
  \[
    \norm{\vec{v}_{1}}
    \pause
    = \sqrt{1^{2}+(-2)^{2}}
    \pause
    = \sqrt{5}
  \]

  \pause
  Autovetor unitário
  \[
    \pause
    \vec{u}_{1}
    \pause
    = \frac{1}{\sqrt{5}}
      \vetor{1}{-2}
  \]
\end{resolution}

%------------------------------------------------------------------------------%
\begin{resolution}{Autovetores para $\lambda=6$}
\begin{columns}
\column{0.4\textwidth}
  \begin{align*}
    \onslide<+->{(A-6I)\vec{x}&=\vec{0}}
    \\
    \onslide<+->{\begin{bmatrix}
      -1 & 2  \\
      2  & -4
    \end{bmatrix}
    \vetor{x}{y} & = \vetor{0}{0}}
  \end{align*}
  \begin{align*}
    \onslide<+->{-x + 2y & = 0 \\}
    \onslide<+->{ x      & = 2y}
  \end{align*}
\column{0.4\textwidth}
  \onslide<+->{Escolhendo $y=1$}
  \onslide<+->{\[
    x = 2
  \]}
  \onslide<+->{\[
    \vec{v}_2 = \vetor{2}{1}
  \]}
  \onslide<+->{\[
    E(6) = \ger \left\{ \vetor{2}{1} \right\}
  \]}
\end{columns}
\end{resolution}

%------------------------------------------------------------------------------%
\begin{resolution}{Normalização de $\vec{v}_2$}
  Norma de
  \(
    \vec{v}_2 = \vetor{2}{1}
  \)

  \pause
  \[
    \norm{\vec{v}_{2}}
    \pause
    = \sqrt{2^{2}+1^{2}}
    \pause
    = \sqrt{5}
  \]

  \pause
  Autovetor unitário
  \[
    \pause
    \vec{u}_{2}
    \pause
    = \frac{1}{\sqrt{5}}
      \vetor{2}{1}
  \]
\end{resolution}

%------------------------------------------------------------------------------%
\begin{resolution}{Diagonalização}
  Organizamos os autovetores unitário nas colunas de $P$
  \[
    P =
    \frac{1}{\sqrt{5}}
    \begin{pmatrix}
      1  & 2 \\
      -2 & 1
    \end{pmatrix}
  \]

  \pause
  A matriz diagonal correspondente é
  \[
    D =
    \begin{pmatrix}
        1 & 0\\
        0 & 6
    \end{pmatrix}
  \]
\end{resolution}

%------------------------------------------------------------------------------%

%%------------------------------------------------------------------------------%
\begin{question}
    Diagonalize ortogonalmente a matriz
    \[
      B =
      \begin{pmatrix}
        4 & 0 \\
        0 & 4
      \end{pmatrix}
    \]
    Determine as matrizes $P$ e $D$ tais que
    \(
      B = P D P^t
    \)

    \bigskip
    \textit{Essa matriz já é diagonal, este exercício é apenas um treino}
\end{question}

%------------------------------------------------------------------------------%
\begin{resolution}{Autovalores}
  Polinômio característico
  \pause
  \[
    p(\lambda)
    \pause
    = \det(B-\lambda I)
    \pause
    = \begin{vmatrix}
      4-\lambda & 0         \\
      0         & 4-\lambda
    \end{vmatrix}
    \pause
    = (4-\lambda)^{2}
  \]

  \pause
  Autovalor
  \[
    \lambda = 4
  \]
\end{resolution}

%------------------------------------------------------------------------------%
\begin{resolution}{Autoespaço}
\begin{columns}[t]
\column{0.4\textwidth}
  \onslide<+->{Equação característica}
  \begin{align*}
    \onslide<+->{(B-4I)\vec{x}&=\vec{0}}
    \\
    \onslide<+->{\begin{pmatrix}
      0 & 0 \\
      0 & 0
    \end{pmatrix}
    \begin{pmatrix}
        x\\y
    \end{pmatrix}
    &=
    \begin{pmatrix}
        0\\0
    \end{pmatrix}}
  \end{align*}
  \onslide<+->{Não há restrições}
\column{0.4\textwidth}
  \onslide<+->{Autoespaço
  \[
    E(4) = \R^2
  \]}
\end{columns}
\end{resolution}

%------------------------------------------------------------------------------%
\begin{resolution}{Base ortonormal}
  Podemos escolher os vetores ortonormais
  \[
    \pause
    \vec{u}_{1} =
    \begin{pmatrix}
      1 \\
      0
    \end{pmatrix}
    \qquad\qquad
    \pause
    \vec{u}_{2} =
    \begin{pmatrix}
      0 \\
      1
    \end{pmatrix}
  \]
  \pause
  \begin{align*}
    \norm{\vec{u}_{1}}          & =1 \\
    \norm{\vec{u}_{2}}          & =1 \\
    \vec{u}_{1}\cdot\vec{u}_{2} & =0
  \end{align*}
\end{resolution}

%------------------------------------------------------------------------------%
\begin{resolution}{Diagonalização}
  Organizamos os autovetores nas colunas de $P$
  \[
    P =
    \begin{pmatrix}
      1 & 0 \\
      0 & 1
    \end{pmatrix}
  \]

  A matriz diagonal é
  \[
    D =
    \begin{pmatrix}
      4 & 0 \\
      0 & 4
    \end{pmatrix}
  \]

  Portanto
  \[
    B = P D P^t
  \]
\end{resolution}

%------------------------------------------------------------------------------%

%%------------------------------------------------------------------------------%
\begin{question}
  Diagonalize ortogonalmente a matriz
  \[
    C =
    \begin{pmatrix}
      1 & 2 \\
      2 & 1
    \end{pmatrix}
  \]
  Determine as matrizes $P$ e $D$ tais que
  \(
    C = P D P^t
  \)
\end{question}

%------------------------------------------------------------------------------%
\begin{resolution}{Autovalores}
\begin{columns}
\column{0.4\textwidth}
  \onslide<+->{Polinômio característico}
  \begin{align*}
    \onslide<+->{p(\lambda)}
    \onslide<+->{& = \det(C-\lambda I)}
    \\
    \onslide<+->{& = \begin{vmatrix}[c]
      1-\lambda & 2         \\
      2         & 1-\lambda
    \end{vmatrix}}
    \\
    \onslide<+->{& = (1-\lambda)^{2} - 4 \\}
    \onslide<+->{& = \lambda^{2} - 2\lambda - 3}
  \end{align*}
\column{0.6\textwidth}
  \[
    \onslide<+->{\Delta}
    \onslide<+->{= b^2 - 4ac}
    \onslide<+->{= (-2)^2 - 4 \times 1 \times (-3)}
    \onslide<+->{= 16}
  \]

  \smallskip
  \[
    \onslide<+->{\sqrt{\Delta}}
    \onslide<+->{= \sqrt{16}}
    \onslide<+->{= 4}
  \]

  \smallskip
  \[
    \onslide<+->{\lambda_1}
    \onslide<+->{= \frac{-b-\sqrt{\Delta}}{2a}}
    \onslide<+->{= \frac{2-4}{2}}
    \onslide<+->{= -1}
  \]

  \smallskip
  \[
    \onslide<+->{\lambda_2}
    \onslide<+->{= \frac{-b+\sqrt{\Delta}}{2a}}
    \onslide<+->{= \frac{2+4}{2}}
    \onslide<+->{= 3}
  \]
\end{columns}
\end{resolution}

%------------------------------------------------------------------------------%
\begin{resolution}{Autovetores para $\lambda=-1$}
\begin{columns}
\column{0.4\textwidth}
  \begin{align*}
    \onslide<+->{(C-I)\vec{x}&=\vec{0}}
    \\
    \onslide<+->{\begin{bmatrix}
      2 & 2 \\
      2 & 2
    \end{bmatrix}
    \vetor{x}{y} & = \vetor{0}{0}}
  \end{align*}
  \begin{align*}
    \onslide<+->{2x + 2y = 0} \\
    \onslide<+->{y = -x}
  \end{align*}
\column{0.4\textwidth}
  \onslide<+->{Escolhendo $x=1$}
  \onslide<+->{\[
    y = -1
  \]}
  \onslide<+->{\[
    \vec{v}_1 = \vetor{1}{-1}
  \]}
  \onslide<+->{\[
    E(1) = \ger \left\{ \vetor{1}{-1} \right\}
  \]}
\end{columns}
\end{resolution}

%------------------------------------------------------------------------------%
\begin{resolution}{Normalização de $\vec{v}_1$}
  Norma de
  \(
    \vec{v}_1 = \vetor{1}{-1}
  \)
  \pause
  \[
    \norm{\vec{v}_{1}}
    \pause
    = \sqrt{1^2 + (-1)^2}
    \pause
    = \sqrt{2}
  \]

  \pause
  Autovetor unitário
  \[
    \pause
    \vec{u}_{1}
    \pause
    = \frac{1}{\sqrt{2}}
      \vetor{1}{-1}
  \]
\end{resolution}

%------------------------------------------------------------------------------%
\begin{resolution}{Autovetores para $\lambda=3$}
\begin{columns}
\column{0.4\textwidth}
  \begin{align*}
    \onslide<+->{(C-I)\vec{x}&=\vec{0}}
    \\
    \onslide<+->{\begin{bmatrix}
      -2 & 2  \\
      2  & -2
    \end{bmatrix}
    \vetor{x}{y} & = \vetor{0}{0}}
  \end{align*}
  \begin{align*}
    \onslide<+->{2x - 2y = 0} \\
    \onslide<+->{y = x}
  \end{align*}
\column{0.4\textwidth}
  \onslide<+->{Escolhendo $x=1$}
  \onslide<+->{\[
    y = 1
  \]}
  \onslide<+->{\[
    \vec{v}_2 = \vetor{1}{1}
  \]}
  \onslide<+->{\[
    E(3) = \ger \left\{ \vetor{1}{1} \right\}
  \]}
\end{columns}
\end{resolution}


%------------------------------------------------------------------------------%
\begin{resolution}{Normalização de $\vec{v}_2$}
  Norma de
  \(
    \vec{v}_2 = \vetor{1}{1}
  \)
  \pause
  \[
    \norm{\vec{v}_{2}}
    \pause
    = \sqrt{1^2 + 1^2}
    \pause
    = \sqrt{2}
  \]

  \pause
  Autovetor unitário
  \[
    \pause
    \vec{u}_{2}
    \pause
    = \frac{1}{\sqrt{2}}
      \vetor{1}{1}
  \]
\end{resolution}

%------------------------------------------------------------------------------%
\begin{resolution}{Diagonalização}
  Organizamos os autovetores nas colunas de $P$
  \[
    P =
    \frac{1}{\sqrt{2}}
    \begin{pmatrix}
      1  & 1 \\
      -1 & 1
    \end{pmatrix}
  \]

  \pause
  A matriz diagonal é
  \[
    D =
    \begin{pmatrix}
      -1 & 0 \\
      0  & 3
    \end{pmatrix}
  \]

  \pause
  Portanto
  \[
    C = P D P^t
  \]
\end{resolution}

%------------------------------------------------------------------------------%

%------------------------------------------------------------------------------%
\begin{question}
  Diagonalize ortogonalmente a matriz
  \[
    A =
    \begin{pmatrix}
      2 & 1 & 0 \\
      1 & 2 & 0 \\
      0 & 0 & 4
    \end{pmatrix}
  \]
  Determine as matrizes $P$ e $D$ tais que
  \(
    A = P D P^t
  \)
\end{question}

%------------------------------------------------------------------------------%
\begin{resolution}{Autovalores}
  Polinômio característico
  \begin{align*}
    \onslide<+->{p(\lambda)}
    \onslide<+->{= \det(A-\lambda I) }
    \onslide<+->{& = \begin{vmatrix}
      2-\lambda & 1         & 0         \\
      1         & 2-\lambda & 0         \\
      0         & 0         & 4-\lambda
    \end{vmatrix}}
    \\
    \onslide<+->{& = (-1)^{3+3} (4-\lambda)
      \begin{vmatrix}
        2-\lambda & 1         \\
        1         & 2-\lambda
      \end{vmatrix}}
    \\
    \onslide<+->{& = (4-\lambda) \left[ (2-\lambda)^2 - 1            \right]} \\
    \onslide<+->{& = (4-\lambda) \left[ 4 - 4\lambda + \lambda^2 - 1 \right]} \\
    \onslide<+->{& = (4-\lambda) \left( \lambda^2 - 4\lambda + 3     \right)}
  \end{align*}
\end{resolution}

%------------------------------------------------------------------------------%
\begin{resolution}{Autovalores}
  \[
    \onslide<+->{(4-\lambda) \left( \lambda^2 - 4\lambda + 3 \right) = 0}
  \]

  \[
    \onslide<+->{\lambda_1 = 4}
    \onslide<+->{\qquad\text{ ou }\qquad
    \lambda^2 - 4\lambda + 3 = 0}
  \]

  \[
    \onslide<+->{\Delta}
    \onslide<+->{= b^2 - 4ac}
    \onslide<+->{= (-4)^2 - 4 \times 1 \times 3}
    \onslide<+->{= 4}
    \qquad\qquad
    \onslide<+->{\sqrt{\Delta}}
    \onslide<+->{= \sqrt{4}}
    \onslide<+->{= 2}
  \]

  \[
    \onslide<+->{\lambda_2}
    \onslide<+->{= \frac{-b+\sqrt{\Delta}}{2a}}
    \onslide<+->{= \frac{4+2}{2}}
    \onslide<+->{= 3}
  \]

  \[
    \onslide<+->{\lambda_3}
    \onslide<+->{= \frac{-b-\sqrt{\Delta}}{2a}}
    \onslide<+->{= \frac{4-2}{2}}
    \onslide<+->{= 1}
  \]
\end{resolution}

%------------------------------------------------------------------------------%
\begin{resolution}{Autovetores para $\lambda_1=4$}
\begin{columns}
\column{0.4\textwidth}
  \begin{align*}
    \onslide<+->{(A-4I)\vec{x}&=\vec{0}}
    \\
    \onslide<+->{\begin{bmatrix}
      -2 & 1  & 0 \\
      1  & -2 & 0 \\
      0  & 0  & 0
    \end{bmatrix}
    \vetortri{x}{y}{z} & = \vetortri{0}{0}{0}}
  \end{align*}
  \vspace{-\baselineskip}
  \begin{align*}
    \onslide<+->{-2x +  y & = 0 \\}
    \onslide<+->{y & = 2x}
    \\[2mm]
    \onslide<+->{x - 2y & =}
    \onslide<+->{x - 4x =}
    \onslide<+->{-x = 0} \\
    \onslide<+->{y & = 0}
  \end{align*}
\column{0.4\textwidth}
  \onslide<+->{Escolhendo $z=1$}
  \onslide<+->{\[
    \vec{v}_1 = \vetortri{0}{0}{1}
  \]}
  \onslide<+->{\[
    E(4) = \ger \left\{ \vetortri{0}{0}{1} \right\}
  \]}
\end{columns}
\end{resolution}

%------------------------------------------------------------------------------%
\begin{resolution}{Normalização de $\vec{v}_1$}
  Norma de
  \(
    \vec{v}_1 = \vetortri{0}{0}{1}
  \)
  \pause
  \[
    \norm{\vec{v}_{2}}
    \pause
    = \sqrt{0^2 +0^2 + 1^2}
    \pause
    = 1
  \]

  \pause
  Autovetor unitário
  \[
    \pause
    \vec{u}_1
    \pause
    = \vetortri{0}{0}{1}
  \]
\end{resolution}

%------------------------------------------------------------------------------%
\begin{resolution}{Autovetores para $\lambda_2=3$}
\begin{columns}
\column{0.4\textwidth}
  \begin{align*}
    \onslide<+->{(A-3I)\vec{x}&=\vec{0}}
    \\
    \onslide<+->{\begin{bmatrix}
      -1 & 1  & 0 \\
      1  & -1 & 0 \\
      0  & 0  & 1
    \end{bmatrix}
    \vetortri{x}{y}{z} & = \vetortri{0}{0}{0}}
  \end{align*}
  \onslide<+->{\begin{align*}
    -x + y & = 0 \\
    z      & = 0
  \end{align*}}
\column{0.4\textwidth}
  \onslide<+->{Escolhendo $x=1$}
  \onslide<+->{\[
    y = 1
  \]}
  \onslide<+->{\[
    \vec{v}_2 = \vetortri{1}{1}{0}
  \]}
  \onslide<+->{\[
    E(3) = \ger \left\{ \vetortri{1}{1}{0} \right\}
  \]}
\end{columns}
\end{resolution}

%------------------------------------------------------------------------------%
\begin{resolution}{Normalização de $\vec{v}_2$}
  Norma de
  \(
    \vec{v}_2 = \vetortri{1}{1}{0}
  \)
  \pause
  \[
    \norm{\vec{v}_{2}}
    \pause
    = \sqrt{1^2 +1^2 + 0^2}
    \pause
    = \sqrt{2}
  \]

  \pause
  Autovetor unitário
  \[
    \pause
    \vec{u}_2
    \pause
    = \frac{1}{\sqrt{2}}
      \vetortri{1}{1}{0}
  \]
\end{resolution}

%------------------------------------------------------------------------------%
\begin{resolution}{Autovetores para $\lambda_3=1$}
\begin{columns}
\column{0.4\textwidth}
  \begin{align*}
    \onslide<+->{(A-I)\vec{x}&=\vec{0}}
    \\
    \onslide<+->{\begin{bmatrix}
      1 & 1 & 0 \\
      1 & 1 & 0 \\
      0 & 0 & 3
    \end{bmatrix}
    \vetortri{x}{y}{z} & = \vetortri{0}{0}{0}}
  \end{align*}
  \onslide<+->{\begin{align*}
    x+y & = 0 \\
    z   & = 0
  \end{align*}}
\column{0.4\textwidth}
  \onslide<+->{Escolhendo $x=1$}
  \onslide<+->{\[
    y = -1
  \]}
  \onslide<+->{\[
    \vec{v}_3 = \vetortri{1}{-1}{0}
  \]}
  \onslide<+->{\[
    E(1) = \ger \left\{ \vetortri{1}{-1}{0} \right\}
  \]}
\end{columns}
\end{resolution}

%------------------------------------------------------------------------------%
\begin{resolution}{Normalização de $\vec{v}_3$}
  Norma de
  \(
    \vec{v}_3 = \vetortri{1}{-1}{0}
  \)
  \pause
  \[
    \norm{\vec{v}_{3}}
    \pause
    = \sqrt{1^2 +(-1)^2 + 0^2}
    \pause
    = \sqrt{2}
  \]

  \pause
  Autovetor unitário
  \[
    \pause
    \vec{u}_3
    \pause
    = \frac{1}{\sqrt{2}}
      \vetortri{1}{-1}{0}
  \]
\end{resolution}

%------------------------------------------------------------------------------%
\begin{resolution}{Base Ortonormal}
  Autovetores unitários
  \[
    \vec{u}_{1} = \vetortri{0}{0}{1}                     \qquad
    \vec{u}_{2} = \frac{1}{\sqrt{2}} \vetortri{1}{1}{0}  \qquad
    \vec{u}_{3} = \frac{1}{\sqrt{2}} \vetortri{1}{-1}{0}
  \]
\end{resolution}

%------------------------------------------------------------------------------%
\begin{resolution}{Diagonalização}
  Organizamos os autovetores nas colunas de $P$
  \[
    P =
    \begin{pmatrix}
      \dfrac{1}{\sqrt{2}}  & \dfrac{1}{\sqrt{2}} & 0 \\[0.3cm]
      -\dfrac{1}{\sqrt{2}} & \dfrac{1}{\sqrt{2}} & 0 \\[0.3cm]
      0                    & 0                   & 1
    \end{pmatrix}
  \]

  Organizamos os autovalores na diagonal de $D$
  \[
    D =
    \begin{pmatrix}
      1 & 0 & 0 \\
      0 & 3 & 0 \\
      0 & 0 & 4
    \end{pmatrix}
  \]
\end{resolution}

%------------------------------------------------------------------------------%

%------------------------------------------------------------------------------%
\begin{question}
  Sabendo que os autovalores da matriz
  \[
    B =
    \begin{pmatrix}
      3 & 1 & 1 \\
      1 & 3 & 1 \\
      1 & 1 & 3
    \end{pmatrix}
  \]
  são
  \(
    \lambda_{1}=5
  \)
  e
  \(
    \lambda_{2}=2
  \)

  \bigskip
  Determine as matrizes $P$ e $D$ tais que
  \(
    B = P D P^t
  \)
\end{question}

%------------------------------------------------------------------------------%
\begin{resolution}{Autovalores}
  Calculamos o polinômio característico
  \begin{align*}
      p(\lambda)
      &=\det(B-\lambda I)\\
      &=
      \begin{vmatrix}
          3-\lambda & 1 & 1\\
          1 & 3-\lambda & 1\\
          1 & 1 & 3-\lambda
      \end{vmatrix}\\
      &=(5-\lambda)(2-\lambda)^2
  \end{align*}

  Portanto
  \[
      \lambda_{1}=5
      \qquad
      \lambda_{2}=2
  \]

  O autovalor $2$ possui multiplicidade algébrica dois
\end{resolution}

%------------------------------------------------------------------------------%
\begin{resolution}{Autoespaço para $\lambda=5$}
  Resolvemos
  \begin{align*}
      (B-5I)\vec{x}&=\vec{0}\\
      \begin{pmatrix}
          -2 & 1 & 1\\
          1 & -2 & 1\\
          1 & 1 & -2
      \end{pmatrix}
      \vetortri{x}{y}{z}
      &=\vetortri{0}{0}{0}
  \end{align*}

  O sistema equivalente fornece
  \begin{align*}
      -2x+y+z&=0\\
      x-2y+z&=0\\
      x+y-2z&=0
  \end{align*}

  Subtraindo as equações
  \begin{align*}
      x&=y\\
      y&=z
  \end{align*}
\end{resolution}

%------------------------------------------------------------------------------%
\begin{resolution}{Autoespaço para $\lambda=5$}
  Portanto
  \[
      E_{5}=
      \operatorname{span}
      \left\{
          \vetortri{1}{1}{1}
      \right\}
  \]

  Escolhemos
  \[
      \vec{v}_{1}=\vetortri{1}{1}{1}
  \]

  Sua norma é
  \[
      \norm{\vec{v}_{1}}=\sqrt{3}
  \]

  O autovetor unitário é
  \[
      \vec{u}_{1}=
      \frac{1}{\sqrt{3}}
      \vetortri{1}{1}{1}
  \]
\end{resolution}

%------------------------------------------------------------------------------%
\begin{resolution}{Autoespaço para $\lambda=2$}
  Resolvemos
  \begin{align*}
      (B-2I)\vec{x}&=\vec{0}\\
      \begin{pmatrix}
          1 & 1 & 1\\
          1 & 1 & 1\\
          1 & 1 & 1
      \end{pmatrix}
      \vetortri{x}{y}{z}
      &=\vetortri{0}{0}{0}
  \end{align*}

  O sistema equivalente é
  \[
      x+y+z=0
  \]

  Escolhendo $x=s$ e $y=t$
  \[
      z=-s-t
  \]
\end{resolution}

%------------------------------------------------------------------------------%
\begin{resolution}{Autoespaço para $\lambda=2$}
  Portanto
  \begin{align*}
      E_{2}
      &=
      \left\{
          \vetortri{s}{t}{-s-t}
          \mid s,t\in\mathbb{R}
      \right\}\\
      &=
      \operatorname{span}
      \left\{
          \vetortri{1}{0}{-1}
          ,
          \vetortri{0}{1}{-1}
      \right\}
  \end{align*}

  Temos
  \[
      \dim(E_{2})=2
  \]
\end{resolution}

%------------------------------------------------------------------------------%
\begin{resolution}{Ortogonalização}
  Escolhemos os vetores
  \[
      \vec{v}_{2}=\vetortri{1}{0}{-1}
      \qquad
      \vec{v}_{3}=\vetortri{0}{1}{-1}
  \]

  Calculamos
  \begin{align*}
      \vec{v}_{2}\cdot\vec{v}_{3}
      &=-1\\
      \norm{\vec{v}_{2}}^{2}
      &=2
  \end{align*}

  Aplicamos Gram-Schmidt
  \begin{align*}
      \vec{w}_{2}
      &=\vec{v}_{2}\\
      \vec{w}_{3}
      &=\vec{v}_{3}
      -\frac{\vec{v}_{3}\cdot\vec{w}_{2}}
      {\norm{\vec{w}_{2}}^{2}}\vec{w}_{2}\\
      &=\vec{v}_{3}+\frac{1}{2}\vec{v}_{2}
  \end{align*}
\end{resolution}

%------------------------------------------------------------------------------%
\begin{resolution}{Ortogonalização}
  Substituindo
  \begin{align*}
      \vec{w}_{3}
      &=
      \vetortri{0}{1}{-1}
      +\frac{1}{2}\vetortri{1}{0}{-1}\\
      &=
      \vetortri{\frac{1}{2}}{1}{-\frac{3}{2}}
  \end{align*}

  Podemos multiplicar por $2$ sem alterar a direção
  \[
      \vec{w}_{3}=\vetortri{1}{2}{-3}
  \]

  Verificamos
  \begin{align*}
      \vec{w}_{2}\cdot\vec{w}_{3}
      &=1-3=0
  \end{align*}
\end{resolution}

%------------------------------------------------------------------------------%
\begin{resolution}{Normalização}
  Calculamos as normas
  \begin{align*}
      \norm{\vec{w}_{2}}
      &=\sqrt{1^{2}+0^{2}+(-1)^{2}}\\
      &=\sqrt{2}\\
      \norm{\vec{w}_{3}}
      &=\sqrt{1^{2}+2^{2}+(-3)^{2}}\\
      &=\sqrt{14}
  \end{align*}

  Os vetores unitários são
  \begin{align*}
      \vec{u}_{2}
      &=\frac{1}{\sqrt{2}}
      \vetortri{1}{0}{-1}\\
      \vec{u}_{3}
      &=\frac{1}{\sqrt{14}}
      \vetortri{1}{2}{-3}
  \end{align*}
\end{resolution}

%------------------------------------------------------------------------------%
\begin{resolution}{Base Ortonormal}
  Obtivemos a base ortonormal de autovetores
  \begin{align*}
      \vec{u}_{1}
      &=\frac{1}{\sqrt{3}}
      \vetortri{1}{1}{1}\\
      \vec{u}_{2}
      &=\frac{1}{\sqrt{2}}
      \vetortri{1}{0}{-1}\\
      \vec{u}_{3}
      &=\frac{1}{\sqrt{14}}
      \vetortri{1}{2}{-3}
  \end{align*}

  O primeiro vetor pertence a $E_{5}$

  Os outros dois pertencem a $E_{2}$
\end{resolution}

%------------------------------------------------------------------------------%
\begin{resolution}{Matriz ortogonal}
  Organizamos os autovetores nas colunas de $P$
  \[
      P=
      \begin{pmatrix}
          \frac{1}{\sqrt{3}} &
          \frac{1}{\sqrt{2}} &
          \frac{1}{\sqrt{14}}\\[0.2cm]
          \frac{1}{\sqrt{3}} &
          0 &
          \frac{2}{\sqrt{14}}\\[0.2cm]
          \frac{1}{\sqrt{3}} &
          -\frac{1}{\sqrt{2}} &
          -\frac{3}{\sqrt{14}}
      \end{pmatrix}
  \]

  Como as colunas formam uma base ortonormal
  \[
      P^tP=I
  \]
\end{resolution}

%------------------------------------------------------------------------------%
\begin{resolution}{Matriz diagonal}
  Os autovalores correspondentes aos vetores das colunas de $P$ são
  \[
      5\qquad 2\qquad 2
  \]

  Portanto
  \[
      D=
      \begin{pmatrix}
          5 & 0 & 0\\
          0 & 2 & 0\\
          0 & 0 & 2
      \end{pmatrix}
  \]

  Assim
  \[
      B=PDP^t
  \]
\end{resolution}

%------------------------------------------------------------------------------%
\begin{resolution}{Resposta}
  A matriz ortogonal é
  \[
      P=
      \begin{pmatrix}
          \frac{1}{\sqrt{3}} &
          \frac{1}{\sqrt{2}} &
          \frac{1}{\sqrt{14}}\\
          \frac{1}{\sqrt{3}} &
          0 &
          \frac{2}{\sqrt{14}}\\
          \frac{1}{\sqrt{3}} &
          -\frac{1}{\sqrt{2}} &
          -\frac{3}{\sqrt{14}}
      \end{pmatrix}
  \]

  A matriz diagonal é
  \[
      D=
      \begin{pmatrix}
          5 & 0 & 0\\
          0 & 2 & 0\\
          0 & 0 & 2
      \end{pmatrix}
  \]

  Portanto
  \[
      B=PDP^t
  \]
\end{resolution}

%------------------------------------------------------------------------------%


%%------------------------------------------------------------------------------%
%\section{Lista Mínima}
%
%%------------------------------------------------------------------------------%
%\begin{frame}{Lista Mínima}
%  \begin{enumerate}
%    \item
%  \end{enumerate}
%
%  \vfill
%  Atenção: A prova é baseada no livro, não nas apresentações
%\end{frame}

%------------------------------------------------------------------------------%
\end{document}
%------------------------------------------------------------------------------%
